\documentclass[varwidth,border=10pt]{standalone}

% ---------- Encoding & Language ----------
\usepackage[T1]{fontenc}
\usepackage[utf8]{inputenc}
\usepackage[ngerman]{babel}

% ---------- Font ----------
\usepackage{sourcesanspro}
\renewcommand{\familydefault}{\sfdefault}

% ---------- Layout & Typography ----------
\usepackage{microtype}
\usepackage{enumitem}
\usepackage{tabularx}
\usepackage{array}

\setlength{\parindent}{0pt}
\setlist[itemize]{leftmargin=*, itemsep=2pt, topsep=4pt}

\begin{document}
    
    {\LARGE\bfseries Imperativ im Deutschen (Befehle, Aufforderungen)}\par\vspace{6pt}
    
    Der \textbf{Imperativ} wird benutzt, um jemanden direkt zu etwas aufzufordern:
    \textit{Komm!} \textit{Hören Sie zu!}
    
    \vspace{8pt}
    
    
    \section*{1) Imperativ mit \textit{du} (informell, 1 Person)}
    \textbf{Regel:} Verbstamm (meist ohne \textit{-st}); \textit{du} steht normalerweise
    \textbf{nicht} im Satz.\par
    \textbf{Beispiele:}
    \begin{itemize}
        \item \textit{Komm!}
        \item \textit{Lern Deutsch!}
        \item \textit{Schreib mir!}
    \end{itemize}
    
    \textbf{Hinweis:} Bei manchen Verben ist ein \textit{-e} möglich (oft etwas formeller):\par
    \textit{Hör zu!} / \textit{Höre zu!}
    
    \vspace{6pt}
    
    
    \section*{2) Imperativ mit \textit{ihr} (informell, mehrere Personen)}
    \textbf{Regel:} Wie die Präsensform von \textit{ihr} (ohne das Wort \textit{ihr}).\par
    \textbf{Beispiele:}
    \begin{itemize}
        \item \textit{Kommt!}
        \item \textit{Hört zu!}
        \item \textit{Lernt Deutsch!}
    \end{itemize}
    
    \vspace{6pt}
    
    
    \section*{3) Imperativ mit \textit{Sie} (formell)}
    \textbf{Regel:} Verb an Position 1 + \textit{Sie} (meist am Ende). Oft mit
    \textit{bitte}.\par
    \textbf{Beispiele:}
    \begin{itemize}
        \item \textit{Kommen Sie bitte!}
        \item \textit{Hören Sie zu!}
        \item \textit{Füllen Sie das Formular aus!}
    \end{itemize}
    
    \vspace{6pt}
    
    
    \section*{4) Trennbare Verben (z.\,B. \textit{aufstehen}, \textit{zumachen}, \textit{anrufen})}
    \textbf{Regel:} Das Präfix steht \textbf{am Satzende}.\par
    \textbf{Beispiele:}
    \begin{itemize}
        \item \textit{Steh auf!}
        \item \textit{Mach die Tür zu!}
        \item \textit{Ruf mich an!}
    \end{itemize}
    
    \vspace{6pt}
    
    
    \section*{5) Verneinter Imperativ (Negation)}
    \textbf{Regel:} \textit{nicht} oder \textit{kein}.\par
    \textbf{Beispiele:}
    \begin{itemize}
        \item \textit{Komm nicht zu spät!}
        \item \textit{Rauch hier nicht!}
        \item \textit{Mach keinen Lärm!}
    \end{itemize}
    
    \vspace{10pt}
    
    
    \section*{Kurzübersicht (Tabelle)}
    \renewcommand{\arraystretch}{1.25}
    \begin{tabularx}{\linewidth}{|>{\raggedright\arraybackslash}p{2.6cm}|X|X|}
\hline
\textbf{Form} & \textbf{Regel} & \textbf{Beispiel} \\
\hline
\textit{du} & Verbstamm (ohne \textit{-st}), \textit{du} meist weglassen
& \textit{Komm! / Lern! / Schreib!} \\
\hline
\textit{ihr} & wie Präsens-\textit{ihr} (ohne \textit{ihr})
& \textit{Kommt! / Lernt! / Hört zu!} \\
\hline
\textit{Sie} & Verb zuerst + \textit{Sie} (oft mit \textit{bitte})
& \textit{Kommen Sie bitte!} \\
\hline
trennbar & Präfix am Ende
& \textit{Steh auf! / Ruf mich an!} \\
\hline
Negation & \textit{nicht} / \textit{kein}
& \textit{Komm nicht! / Mach keinen Lärm!} \\
\hline
    \end{tabularx}

\end{document}
